% 第一篇：复变函数部分第2至第7章目录规划

\chapter*{第一篇复变函数部分第2至第7章目录规划}
\addcontentsline{toc}{chapter}{第一篇复变函数部分第2至第7章目录规划}

\begin{wulibj}[目录设计的总体思路]
这一份目录规划主要服务于第一篇``复变函数''的整体重排，目标是把主线内容和选学内容分开，让学生在自学时更容易抓住真正的核心推进。

建议把第2章到第5章作为主线必学内容，把第6章和第7章作为选学内容。这样安排以后，整条主线会更清楚：
\[
    \text{解析函数} \Longrightarrow \text{复积分} \Longrightarrow \text{级数展开与奇点} \Longrightarrow \text{留数理论}.
\]

多值函数、支点、解析支和保角映射都很重要，但它们更适合在主线完成以后，再作为选学专题系统展开。这样既能保证理论依赖顺序清楚，也能减少初学者在前半段被过多新概念打断的情况。
\end{wulibj}

\begin{tishi}[使用建议]
\begin{enumerate}
    \item 这一文件是章节规划稿，不是直接替换正文的现成章节。
    \item 每章都只写到``应该讲什么''和``为什么这样安排''，便于后续逐章展开。
    \item 若后续决定继续细化，可在此基础上分别再为每一章补写详细提纲。
\end{enumerate}
\end{tishi}

\section*{第2章\quad 复变函数与解析函数初步}

\subsection*{本章定位}
这一章的任务，是把学生从``会处理复数''带到``理解解析函数为什么是复分析的核心对象''。它是第一篇最关键的桥梁章，必须承担起概念奠基的作用。

\subsection*{建议目录}
\begin{enumerate}
    \item 复平面上的点集、区域与曲线
    \item 复变函数的概念与映射观点
    \item 复变函数的极限与连续
    \item 复变函数的导数与解析性
    \item 柯西--黎曼条件及其应用
    \item 解析函数的调和性质与物理背景
    \item 本章小结
    \item 习题
\end{enumerate}

\subsection*{这样安排的理由}
\begin{itemize}
    \item 第一节先建立区域、连通、单连通、曲线、简单闭曲线等语言，为后面的复积分作准备。
    \item 中间几节围绕``极限--连续--导数--解析''这一主线推进，帮助学生逐步体会复可导比实可导严格得多。
    \item 柯西--黎曼条件既是判别工具，也是学生第一次真正感受到``不同方向趋近必须给出同一结果''的地方。
    \item 最后一节加入调和函数与二维势场背景，可以把复变函数和后面数理方程部分自然接起来。
    \item 本章不正式展开多值函数，只在适当位置提醒：以后会遇到 $\log z$、$z^\alpha$、$\sqrt{z}$ 等不能全局单值处理的函数。
\end{itemize}

\section*{第3章\quad 复积分与柯西理论}

\subsection*{本章定位}
这一章的任务，是让学生第一次真正看到复分析和高等数学中实变积分的根本差别，并建立柯西积分定理与柯西积分公式这两块核心基石。

\subsection*{建议目录}
\begin{enumerate}
    \item 复积分的定义、参数化与基本性质
    \item 曲线积分与路径无关的初步认识
    \item 柯西积分定理
    \item 原函数与复积分的基本公式
    \item 柯西积分公式及其直接推论
    \item 解析函数的若干重要性质
    \item 本章小结
    \item 习题
\end{enumerate}

\subsection*{这样安排的理由}
\begin{itemize}
    \item 先把积分定义、参数化和估值不等式讲清，再进入理论定理，节奏更稳。
    \item ``路径无关''与``原函数存在''是初学者最容易忽略、但又最能体现复积分特殊性的地方，应单独强调。
    \item 柯西积分公式必须作为全章重点，因为后面的泰勒展开、无穷可微性等都从这里生长出来。
    \item 最后一节可整理解析函数的若干重要性质，但不宜铺得过宽，以免冲淡主线。
\end{itemize}

\section*{第4章\quad 解析函数的级数表示与孤立奇点}

\subsection*{本章定位}
这一章要帮助学生建立一种新的眼光：解析函数不仅``可导''，而且在局部有非常整齐的级数结构。也正是在这一章里，学生会第一次系统接触奇点。

\subsection*{建议目录}
\begin{enumerate}
    \item 复数项级数与幂级数预备
    \item 幂级数的收敛半径与基本性质
    \item 解析函数的泰勒展开
    \item 零点与解析函数的局部结构
    \item 洛朗展开
    \item 孤立奇点及其分类
    \item 从局部展开到留数理论
    \item 本章小结
    \item 习题
\end{enumerate}

\subsection*{这样安排的理由}
\begin{itemize}
    \item 预备部分只讲后文真正需要的级数知识，不宜重新铺开一遍高等数学中的全部级数理论。
    \item 泰勒展开应强调：它不是近似技巧，而是解析函数局部结构的真实表达。
    \item 零点、洛朗展开和奇点分类三部分连在一起，正好构成从``局部整齐''走向``局部失效''的自然过渡。
    \item 最后一节专门指出 $(z-z_0)^{-1}$ 项的重要性，顺势把学生带到下一章的留数理论。
\end{itemize}

\section*{第5章\quad 留数理论}

\subsection*{本章定位}
这一章是复变函数部分中计算威力最强的一章。它的核心任务，是让学生掌握如何从洛朗展开中抽取最关键的信息，并把这种局部信息转化成围道积分和实积分的计算工具。

\subsection*{建议目录}
\begin{enumerate}
    \item 留数的定义与常用计算法
    \item 留数定理
    \item 留数定理与实积分计算的一般步骤
    \item 有理型实积分的计算
    \item 三角积分与单位圆方法
    \item 辐角原理与鲁歇定理（选读）
    \item 本章小结
    \item 习题
\end{enumerate}

\subsection*{这样安排的理由}
\begin{itemize}
    \item 本章主线只处理单值解析函数情形下的留数理论，这样理论结构最清楚。
    \item 留数的定义与留数定理必须放在最前面，后面的积分应用都建立在这两节上。
    \item 实积分应用部分宜先讲不涉及支割的典型题型，如有理函数积分和三角积分。
    \item 辐角原理与鲁歇定理虽然重要，但对部分读者来说更偏理论进阶，因此可在章内标为选读。
    \item 若后续需要处理含 $\log z$、$z^\alpha$、$\sqrt{z}$ 的积分，应放到第6章去讲，而不是打断本章主线。
\end{itemize}

\section*{第6章\quad 多值函数、支点与解析支（选学）}

\subsection*{本章定位}
这一章是专门的选学专题。它要解决的问题不是``再介绍几个新函数''，而是回答一个更本质的问题：为什么有些在复分析里非常自然的对象，不能像普通单值函数那样在整个区域上直接处理。

\subsection*{建议目录}
\begin{enumerate}
    \item 多值性从哪里来：从 $\arg z$ 到 $\log z$
    \item 主值、解析支与支割
    \item 支点、绕行现象与多叶直观
    \item 典型多值函数：$\log z$、$z^\alpha$、$\sqrt[n]{z}$
    \item 支割路径与复积分方法
    \item 带多值函数的典型积分举例
    \item 本章小结
    \item 习题
\end{enumerate}

\subsection*{这样安排的理由}
\begin{itemize}
    \item 第1到第4节先把多值函数的概念世界搭起来，避免学生一上来就被技巧性积分淹没。
    \item 第3节中的``多叶''主要用来帮助建立直观，不必把黎曼面理论正式展开成主线。
    \item 第5节是本章与第5章的衔接点，应明确说明：留数定理本身并没有改变，变化的是函数要先选支、路径要避开支割。
    \item 第6节可以安排钥匙孔围道、含 $\log x$、$x^\alpha$、$\sqrt{x}$ 的典型积分，让学生真正看到这一章的用途。
\end{itemize}

\section*{第7章\quad 保角映射初步（选学）}

\subsection*{本章定位}
这一章从解析函数的几何意义出发，把``局部上像旋转加伸缩''这一结论扩展到具体映射的构造与应用中。它既有很强的几何直观，也能和二维静电场、二维流体等问题联系起来。

\subsection*{建议目录}
\begin{enumerate}
    \item 保角映射的概念与局部保角性
    \item 分式线性变换
    \item 典型初等映射
    \item 幂函数与角域映射
    \item 常见区域到标准区域的映射
    \item 保角映射在二维势场中的简单应用
    \item 本章小结
    \item 习题
\end{enumerate}

\subsection*{这样安排的理由}
\begin{itemize}
    \item 第一节先回顾解析函数局部保角的思想，让学生知道后面讨论的不是凭空出现的新主题。
    \item 分式线性变换是最基础、最典型的一类保角映射，必须单独成节。
    \item 幂函数、指数函数等典型映射要强调``区域怎样变化''，而不只是记公式。
    \item 如果第6章已经讲过解析支，那么这里再讲幂函数和根式诱导的角域映射会更顺。
    \item 最后一节应用不宜写得过重，重点是让学生看见保角映射怎样把复杂边界化成标准边界。
\end{itemize}

\section*{建议采用的第一篇后半段目录}

\begin{tishi}[推荐目录总表]
\begin{enumerate}
    \item 第2章\quad 复变函数与解析函数初步
    \item 第3章\quad 复积分与柯西理论
    \item 第4章\quad 解析函数的级数表示与孤立奇点
    \item 第5章\quad 留数理论
    \item 第6章\quad 多值函数、支点与解析支（选学）
    \item 第7章\quad 保角映射初步（选学）
\end{enumerate}
\end{tishi}

\begin{jinshi}[编排时需要特别注意的地方]
\begin{enumerate}
    \item 不要在第2章提前正式展开复对数、幂函数和根式函数的多值理论，否则第6章会失去独立存在的意义。
    \item 不要把带支割的积分路径提前塞进第5章主干，否则学生会在还没学会解析支之前先被技巧绊住。
    \item 不要把第6章写成抽象的黎曼面理论导论；对本讲义来说，重点应始终是``帮助学生处理具体函数与具体积分''。
    \item 第7章应保持几何直观和物理应用导向，不宜写成只给变换公式的技巧汇编。
\end{enumerate}
\end{jinshi}
